% Part: computability
% Chapter: recursive-functions
% Section: notations-pr-functions

\documentclass[../../../include/open-logic-section]{subfiles}

\begin{document}

\olfileid{cmp}{rec}{not}
\olsection{نمادگذاری‌های بازگشت اولیه}

یکی از مزایای داشتن توصیف استقراییِ دقیق از توابع بازگشتی اولیه این
است که می‌توانیم آن‌ها را به‌طور نظام‌مند توصیف کنیم. برای نمونه،
می‌توانیم به هر تابع از این نوع «نمادی» نسبت دهیم. نمادهای $\Zero$،
$\Succ$ و~$\Proj{n}{i}$ را به‌ترتیب برای صفر، جانشین و افکنش‌ها به
کار می‌بریم. اکنون فرض کنید $h$ با ترکیب از تابع $k$موضعیِ~$f$ و
توابع $n$موضعیِ $g_0$، \dots،~$g_{k-1}$ تعریف شده و به این توابع
نمادهای $F$، $G_0$، \dots،~$G_{k-1}$ را نسبت داده باشیم. آنگاه با
نماد تازهٔ $\fn{Comp}_{k,n}$ می‌توان تابع $h$ را با
$\fn{Comp}_{k,n}[F,G_0,\dots,G_{k-1}]$ نشان داد.

برای توابع تعریف‌شده با بازگشت اولیه نیز می‌توانیم از نمادگذاری‌های
همانند استفاده کنیم. فرض کنید تابع $(k+1)$موضعیِ~$h$ با بازگشت اولیه
از تابع $k$موضعیِ~$f$ و تابع $(k+2)$موضعیِ~$g$ تعریف شده و نمادهای
نسبت‌داده‌شده به $f$ و~$g$ به‌ترتیب $F$ و~$G$ باشند. آنگاه نماد
نسبت‌داده‌شده به~$h$ برابر $\fn{Rec}_k[F,G]$ است.

به یاد آورید که تابع جمع با بازگشت اولیه چنین تعریف می‌شود:
\begin{align*}
  \Add(x_0, 0) & = \Proj{1}{0}(x_0) = x_0\\
  \Add(x_0, y+1) & = \Succ(\Proj{3}{2}(x_0, y, \Add(x_0, y))) = \Add(x_0, y) +1
\end{align*}
اینجا $\Proj{1}{0}$ نقش~$f$ را دارد و نقش~$g$ را
$\Succ(\Proj{3}{2}(x_0,y,z))$ بر عهده دارد. نماد نسبت‌داده‌شده به
تابع اخیر $\fn{Comp}_{1,3}[\Succ,\Proj{3}{2}]$ است، زیرا این تابع با
ترکیب از تابع یک‌موضعیِ~$\Succ$ و تابع سه‌موضعیِ~$\Proj{3}{2}$ تعریف
شده است. با این قرارداد می‌توان تابع جمع را چنین نشان داد:
\[
\fn{Rec}_1[\Proj{1}{0},\fn{Comp}_{1,3}[\Succ,\Proj{3}{2}]].
\]
این نمادگذاری‌ها گاه سودمندند؛ برای نمونه، هنگامی که توابع بازگشتی
اولیه را برمی‌شماریم.

\begin{prob}
  نمادگذاری کامل بازگشتی اولیهٔ~$\Mult$ را به دست دهید.
\end{prob}

\end{document}
